\subsection{Reproducible Command Flow}

This appendix collects the exact command-line invocations that generate all checkpoints and evaluation files referenced in the main text. All commands are executed from the project root on Windows PowerShell with the virtual environment activated. The expected output path of each training command is a checkpoint file under the \texttt{checkpoints/} tree, and the expected output path of each evaluation command is a variant-specific JSON file written next to the corresponding checkpoint.

\textbf{Environment setup.} The following commands create and activate a clean Python virtual environment and verify the CUDA availability that is required for reproducing all training runs on a single GPU workstation.
\begin{itemize}
    \item \texttt{python -m venv .venv}
    \item \texttt{.\textbackslash .venv\textbackslash Scripts\textbackslash Activate.ps1}
    \item \texttt{python -c}\\
    \texttt{\phantom{xx}"import torch; print(torch.cuda.is\_available())"}
\end{itemize}

\textbf{Training commands.} The following commands reproduce the five training runs reported in Tables~\ref{tab:awgn-main}, \ref{tab:ablation}, and \ref{tab:rayleigh}. The semantic-loss ablation is controlled by passing \texttt{--beta 0.0} and tagging the run with \texttt{--tag no\_sem}.
\begin{itemize}
    \item \texttt{python train.py}\\
    \texttt{\phantom{xx}--model baseline --channel awgn --epochs 80}
    \item \texttt{python train.py}\\
    \texttt{\phantom{xx}--model improved --channel awgn --epochs 80}
    \item \texttt{python train.py}\\
    \texttt{\phantom{xx}--model improved --channel awgn --epochs 80}\\
    \texttt{\phantom{xx}--beta 0.0 --tag no\_sem}
    \item \texttt{python train.py}\\
    \texttt{\phantom{xx}--model baseline --channel rayleigh --epochs 80}
    \item \texttt{python train.py}\\
    \texttt{\phantom{xx}--model improved --channel rayleigh --epochs 80}
\end{itemize}

\textbf{Evaluation commands.} The following evaluation commands reproduce the quantitative results reported in the main text. The \texttt{--beam\_width} flag switches decoding from greedy to beam search, and the \texttt{--use\_llm} flag together with \texttt{--llm\_model} enables the fine-tuned BERT refiner.
\begin{itemize}
    \item \texttt{python evaluate.py}\\
    \texttt{\phantom{xx}--checkpoint checkpoints/baseline\_awgn/best.pt}\\
    \texttt{\phantom{xx}--model baseline --channel awgn}
    \item \texttt{python evaluate.py}\\
    \texttt{\phantom{xx}--checkpoint checkpoints/improved\_awgn/best.pt}\\
    \texttt{\phantom{xx}--model improved --channel awgn}
    \item \texttt{python evaluate.py}\\
    \texttt{\phantom{xx}--checkpoint}\\
    \texttt{\phantom{xx}checkpoints/improved\_awgn\_no\_sem/best.pt}\\
    \texttt{\phantom{xx}--model improved --channel awgn}
    \item \texttt{python evaluate.py}\\
    \texttt{\phantom{xx}--checkpoint checkpoints/baseline\_rayleigh/best.pt}\\
    \texttt{\phantom{xx}--model baseline --channel rayleigh}
    \item \texttt{python evaluate.py}\\
    \texttt{\phantom{xx}--checkpoint checkpoints/improved\_rayleigh/best.pt}\\
    \texttt{\phantom{xx}--model improved --channel rayleigh}
    \item \texttt{python evaluate.py}\\
    \texttt{\phantom{xx}--checkpoint checkpoints/improved\_awgn/best.pt}\\
    \texttt{\phantom{xx}--model improved --channel awgn}\\
    \texttt{\phantom{xx}--beam\_width 5 --eval\_tag beam5}
    \item \texttt{python evaluate.py}\\
    \texttt{\phantom{xx}--checkpoint checkpoints/improved\_awgn/best.pt}\\
    \texttt{\phantom{xx}--model improved --channel awgn --use\_llm}\\
    \texttt{\phantom{xx}--llm\_model checkpoints/bert\_europarl}\\
    \texttt{\phantom{xx}--eval\_tag ftbert}
    \item \texttt{python evaluate.py}\\
    \texttt{\phantom{xx}--checkpoint checkpoints/improved\_awgn/best.pt}\\
    \texttt{\phantom{xx}--model improved --channel awgn --beam\_width 5}\\
    \texttt{\phantom{xx}--use\_llm}\\
    \texttt{\phantom{xx}--llm\_model checkpoints/bert\_europarl}\\
    \texttt{\phantom{xx}--eval\_tag beam5\_ftbert}
\end{itemize}

\subsection{Data Files Used for Quantitative Analysis}

All tables in the main text are computed directly from the following JSON files. Every file is generated by the corresponding evaluation command above, stored under the \texttt{checkpoints/} tree, and is amenable to version control.

\begin{itemize}
    \item \texttt{checkpoints/baseline\_awgn/eval\_awgn.json}
    \item \texttt{checkpoints/improved\_awgn/eval\_awgn.json}
    \item \texttt{checkpoints/improved\_awgn\_no\_sem/}\\
    \texttt{\phantom{xx}eval\_awgn.json}
    \item \texttt{checkpoints/baseline\_rayleigh/eval\_rayleigh.json}
    \item \texttt{checkpoints/improved\_rayleigh/eval\_rayleigh.json}
    \item \texttt{checkpoints/improved\_awgn/eval\_awgn\_beam5.json}
    \item \texttt{checkpoints/improved\_awgn/eval\_awgn\_ftbert.json}
    \item \texttt{checkpoints/improved\_awgn/}\\
    \texttt{\phantom{xx}eval\_awgn\_beam5\_ftbert.json}
\end{itemize}

\subsection{Notes on Figure and Table Preparation}

To satisfy the report template requirements, the following conventions are followed consistently throughout the document. First, every table and figure must be referenced in the main text before its appearance, so that the reader always encounters a pointer before the object itself. Second, captions are written to be informative and self-contained, so that a reader who only glances at the table can still interpret the main message. Third, axis labels and table column headers include units where relevant (for example \si{dB} for SNR), and every SNR sweep explicitly states the corresponding channel model so that readers cannot confuse AWGN and Rayleigh rows by accident.

For the final submission version, it is also recommended to include at least one figure that visualizes the overall experimental workflow (from data preprocessing through model training and evaluation), and one figure that compares key BLEU-4 trends across the major model variants across the full SNR range. Such visual summaries improve readability for examiners, compress the key message of multi-row tables, and make the evidence chain from design choice to quantitative outcome clearer. Future revisions of this thesis could add these figures without any change to the surrounding text, thanks to the figure-placement conventions described above.

\subsection{Reproducibility Check-list}

The following check-list summarises the minimum steps required to independently reproduce the quantitative claims of the report:
\begin{enumerate}
    \item Prepare the Europarl corpus and run the preprocessing script once to generate tokenized shards.
    \item Execute the five training commands listed above; each one writes a \texttt{best.pt} checkpoint under its own checkpoint folder.
    \item Execute the eight evaluation commands listed above; each one writes a JSON file that can be directly loaded to reproduce the corresponding row of the result tables.
    \item Cross-check the reported numbers in Tables 1--4 against the numerical fields of the JSON files, using the provided comparison helper script.
\end{enumerate}

Following this check-list end-to-end takes roughly two to three days on a single mid-range GPU, with the Rayleigh training runs accounting for the largest share of the total wall-clock time. All hyperparameters are documented either inside the training commands or inside the project configuration file, so that no further undocumented parameter is needed to reproduce the reported figures.
